% !TEX root = ../main.tex
% ============================================================
% Results chapter
% First modern draft: 1D Dirichlet diagnostics and 2D FD/PML solver tests.
% ============================================================

\chapter{Results}
\label{chapter:results}

This chapter evaluates whether learned frequency transfer can be used in a
Helmholtz solver. The central conclusion is deliberately more cautious than a
simple claim of neural acceleration. A learned transfer model can predict
high-frequency wavefields with small relative field error, but Krylov methods
are controlled by residuals, not by field error alone. The Helmholtz operator
can amplify spectrally small field errors into large residual components. The
one-dimensional Dirichlet experiments make this mechanism exact and
diagnosable through analytical eigenpairs. The two-dimensional finite
difference/PML experiments then test whether the same lessons survive in the
solver-facing setting used elsewhere in the thesis.

The chapter is organised around a repeated pattern. First, the field prediction
is measured. Second, the corresponding initial residual and FGMRES behaviour
are measured. Third, the discrepancy between these two views is interpreted in
terms of the object being learned: a complete solution field, a residual
restriction, or a correction prolongation. This distinction is the main result
of the chapter.

% ============================================================
\section{One-dimensional Dirichlet results}
\label{sec:results-1d-dirichlet}
% ============================================================

The one-dimensional study isolates the residual-amplification mechanism in a
setting where the spectral analysis is exact. The problem is the homogeneous
Dirichlet Helmholtz equation on a grid with $N=512$ unknowns and no PML. The
frequency pair is
\begin{equation}
    \omega_L = 16,
    \qquad
    \omega_H = 32.
\end{equation}
With the convention
\begin{equation}
    A_\omega = -D_{xx} - \omega^2 I,
\end{equation}
the Dirichlet eigenpairs are known analytically:
\begin{equation}
    \lambda_k(\omega)
    =
    \frac{4}{h^2}
    \sin^2\!\left(\frac{\pi k}{2(N+1)}\right)
    -
    \omega^2,
    \qquad
    v_k(j)
    =
    \sqrt{\frac{2}{N+1}}
    \sin\!\left(\frac{j k \pi}{N+1}\right).
    \label{eq:results-dirichlet-eigenpairs}
\end{equation}
The eigenvectors in~\eqref{eq:results-dirichlet-eigenpairs} are
Euclidean-normalised. This makes the modal coefficients directly comparable
across the field error and residual error.

The learned solution-transfer model in this section is a one-dimensional
U-Net trained to map $u_{16}$ to $u_{32}$. Its input consists of the real and
imaginary parts of $u_{16}$ together with real and imaginary right-hand-side
channels. The loss is the full-grid complex relative $\ell^2$ field loss. The
best validation loss is approximately $5.37\times 10^{-3}$, showing that the
network is a strong field predictor in this controlled setting.

% ------------------------------------------------------------
\subsection{Warm starts: accurate fields do not imply small residuals}
\label{subsec:results-1d-warm-start}
% ------------------------------------------------------------

Table~\ref{tab:results-1d-warm-start} gives the first solver-facing result.
The neural warm start reduces the field error from $1$ for the zero initial
guess to about $6\times 10^{-2}$. However, this does not translate into a
large GMRES improvement. More importantly, the raw residual of the neural
initial guess is worse than the zero start: approximately $11.6\|b\|$ rather
than $\|b\|$.

\begin{table}[htbp]
\centering
\caption{One-dimensional Dirichlet warm-start result for the
$16\to32$ pair. The learned field is accurate in relative field norm, but its
initial residual is larger than the zero-start residual.}
\label{tab:results-1d-warm-start}
\begin{tabular}{lccc}
\toprule
Initial guess & Relative field error & $\|r_0\|_2/\|b\|_2$ & GMRES iterations \\
\midrule
Zero start & $1.000$ & $1.000$ & $16.0$ \\
Raw Dirichlet U-Net & $0.058$--$0.063$ & $\approx 11.6$ & $15.6$ \\
Exact solution & $0$ & $0$ & $1.0$ \\
\bottomrule
\end{tabular}
\end{table}

The explanation is algebraic. Let $u_\star$ be the exact high-frequency
solution and let $x_0$ be the initial guess. The initial error and residual are
\begin{equation}
    e_0 = u_\star - x_0,
    \qquad
    r_0 = b - A_H x_0 = A_H e_0.
    \label{eq:results-error-residual}
\end{equation}
Expanding in the Dirichlet eigenbasis gives
\begin{equation}
    e_0 = \sum_{k=1}^{N} c_k(e_0)v_k,
    \qquad
    r_0 = \sum_{k=1}^{N} \lambda_k c_k(e_0)v_k.
    \label{eq:results-modal-amplification}
\end{equation}
Thus
\begin{equation}
    c_k(r_0) = \lambda_k c_k(e_0).
\end{equation}
Small field-error coefficients in modes with large $|\lambda_k|$ can therefore
become large residual coefficients. Field accuracy and solver accuracy are not
the same metric.

\begin{figure}[htbp]
    \centering
    \includegraphics[width=\linewidth]{figures/ch6/10_scaled_complex_spectrum.png}
    \caption{Scaled spectrum for the 1D Dirichlet operator and
    the CSL-preconditioned analytical operator. Source plot:
    \texttt{deep\_spectral/10\_scaled\_complex\_spectrum.png}. The
    high-frequency Dirichlet operator has 10 negative modes, 502 positive
    modes, and an absolute condition number of approximately $2.82\times
    10^{4}$, while the CSL-preconditioned analytical spectrum with
    $\beta=0.3$ has condition number approximately $8.29$.}
    \label{fig:results-1d-spectrum}
\end{figure}

The spectral facts quantify the danger. For $A_H$ at $\omega_H=32$, there are
10 negative modes and 502 positive modes. The absolute condition number
$\frac{\max_k |\lambda_k|}{\min_k |\lambda_k|}$ is approximately
$2.82\times 10^4$.
After analytical CSL preconditioning with $\beta=0.3$, the corresponding
condition number is approximately $8.29$. This contrast explains why the
question is not whether the U-Net predicts the field well, but whether its
error lies in modes that matter to the preconditioned Krylov process.

\begin{figure}[htbp]
    \centering
    \includegraphics[width=\linewidth]{figures/ch6/11_error_and_residual_modes_sorted.png}
    \caption{Modal field-error coefficients and residual
    coefficients sorted by $|\lambda_k|$. Source plot:
    \texttt{deep\_spectral/11\_error\_and\_residual\_modes\_sorted.png}.
    The learned field error is small in norm, but multiplication by
    $\lambda_k$ amplifies middle- and high-frequency modal components into a
    large residual.}
    \label{fig:results-1d-error-residual-modes}
\end{figure}

The same diagnosis appears in the residual energy bands. The zero start places
most residual energy in the near-resonant and low-$|\lambda|$ bands. The raw
U-Net removes much of this low-mode content but moves residual energy into the
middle and high-$|\lambda|$ bands. In particular, the model residual has
approximately $0.4926$ of its residual spectral energy in the middle band and
$0.4521$ in the high band, compared with essentially no high-band residual
energy for the zero start. The model is therefore useful, but not uniformly
safe.

\begin{figure}[htbp]
    \centering
    \includegraphics[width=\linewidth]{figures/ch6/13_residual_energy_bands.png}
    \caption{Residual energy by $|\lambda|$ band. Source plot:
    \texttt{deep\_spectral/13\_residual\_energy\_bands.png}. The raw U-Net
    shifts residual energy away from near-resonant modes and into middle/high
    modes, explaining why its field accuracy does not by itself guarantee a
    better Krylov initial state.}
    \label{fig:results-1d-residual-energy-bands}
\end{figure}

The result of this subsection is therefore negative only in the narrow
warm-start sense. It is positive diagnostically: the U-Net has learned useful
low-frequency-to-high-frequency information, but this information must be
filtered or gated before being trusted by the solver.

% ------------------------------------------------------------
\subsection{Spectral filtering and residual gating}
\label{subsec:results-1d-spectral-filtering}
% ------------------------------------------------------------

The spectral filtering experiment asks whether the useful part of the learned
prediction can be kept while harmful modes are removed. Three variants are
especially informative. The \emph{raw U-Net} uses the prediction directly. The
\emph{low5} filter keeps only the lowest 5\% of modal coefficients ordered by
$|\lambda_k|$. The \emph{residual gate} keeps a learned modal component only
when it reduces the modal residual relative to the zero start.

\begin{table}[htbp]
\centering
\caption{Spectral filtering and residual gating for the 1D Dirichlet
$16\to32$ warm start. The learned model contains useful information, but the
solver benefits only when the spectrally harmful components are removed.}
\label{tab:results-1d-filtering}
\begin{tabular}{lccc}
\toprule
Method & Relative field error & $\|r_0\|_2/\|b\|_2$ & GMRES iterations \\
\midrule
Zero start & $1.000$ & $1.000$ & $16.0$ \\
Raw U-Net & $\approx 0.060$ & $\approx 11.68$ & $15.6$ \\
Low5 filter & $\approx 0.058$ & $\approx 0.83$ & $15.6$ \\
Residual gate & $\approx 0.055$ & $\approx 0.716$ & $15.5$ \\
\bottomrule
\end{tabular}
\end{table}

The residual gate gives the clearest interpretation. It improves the field
error slightly relative to the raw model and reduces the initial residual below
the cold-start value. This confirms that the learned model contains useful
solver-facing information. The improvement in GMRES iterations remains small,
however, so this is not yet a complete solver method. Its value is conceptual:
spectral selectivity is necessary.

\begin{figure}[htbp]
    \centering
    \includegraphics[width=\linewidth]{figures/ch6/20_filtering_summary_bars.png}
    \caption{Summary of field error, initial residual, and
    GMRES iterations for raw and filtered 1D warm starts. Source plot:
    \texttt{spectral\_filtering/20\_filtering\_summary\_bars.png}.}
    \label{fig:results-1d-filtering-bars}
\end{figure}

\begin{figure}[htbp]
    \centering
    \includegraphics[width=\linewidth]{figures/ch6/23_gmres_residual_convergence.png}
    \caption{GMRES residual convergence for the zero start,
    raw U-Net, filtered U-Net, residual-gated U-Net, and exact-solution
    reference. Source plot:
    \texttt{spectral\_filtering/23\_gmres\_residual\_convergence.png}.}
    \label{fig:results-1d-filtering-gmres}
\end{figure}

The consequence is that warm-start evaluation alone is too blunt. A learned
field predictor can be good, bad, or useful-but-dangerous depending on where
its error sits in the spectrum. The residual gate turns this observation into
an operational rule: use learned information only where it improves the
solver's residual metric.

% ------------------------------------------------------------
\subsection{Learned preconditioning inside FGMRES}
\label{subsec:results-1d-learned-preconditioner}
% ------------------------------------------------------------

The final one-dimensional experiment places the learned transfer inside
FGMRES as part of a V-cycle-like preconditioner. This is a different use case
from a one-shot warm start. A Krylov preconditioner acts repeatedly on the
current high-frequency residual and must return a correction. It therefore
requires the correct mathematical objects: residual restriction, a low-frequency
solve, and correction prolongation.

Table~\ref{tab:results-1d-preconditioner} summarises the main outcomes. A
CSL-only preconditioner requires about 16 iterations. Combining an exact
low-frequency solve with gated $T_\uparrow$ reduces this to about 12.65
iterations. Raw learned variants can fail catastrophically, and a
solution-trained $T_\downarrow$ does not work as a residual restriction. A
residual-loss learned down/up pair is safer but still does not improve the
iteration count.

\begin{table}[htbp]
\centering
\caption{Learned preconditioning inside FGMRES for the 1D Dirichlet problem.
The successful variant uses an exact low-frequency solve and a gated
high-frequency prolongation. Solution-transfer networks cannot be freely
reused as residual-transfer networks.}
\label{tab:results-1d-preconditioner}
\begin{tabular}{lc}
\toprule
Preconditioner variant & Mean FGMRES iterations \\
\midrule
CSL-only & $\approx 16.0$ \\
Exact low solve + gated $T_\uparrow$ & $\approx 12.65$ \\
Raw learned variants & Can fail catastrophically \\
Solution-trained $T_\downarrow$ as residual restriction & Does not work \\
Residual-loss learned down/up & Safer, but no iteration improvement \\
Fine-tuned learned restriction $R_\theta$ & Reduces to CSL-only \\
\bottomrule
\end{tabular}
\end{table}

The gating statistics sharpen the interpretation. The fine-tuned learned
restriction $R_\theta$ reached validation loss about $0.939$, but the gate
accepted it in $0\%$ of modes, so the resulting method reduced to CSL-only. By
contrast, the exact-low plus $T_\uparrow$ variant was accepted in only about
$2.8\%$ of modal components, but those components were enough to reduce the
iteration count. The successful learned preconditioner is therefore not a
dense replacement for the classical preconditioner. It is a selective
correction mechanism that helps in a small set of solver-relevant modes.

\begin{figure}[htbp]
    \centering
    \includegraphics[width=\linewidth]{figures/ch6/100_fgmres_learned_preconditioner_convergence.png}
    \caption{FGMRES convergence for CSL-only and learned
    preconditioner variants. Source plot:
    \texttt{fgmres\_learned\_preconditioner/100\_fgmres\_learned\_preconditioner\_convergence.png}.}
    \label{fig:results-1d-preconditioner-convergence}
\end{figure}

\begin{figure}[htbp]
    \centering
    \includegraphics[width=\linewidth]{figures/ch6/120_residual_improvement_ratio_vs_csl.png}
    \caption{Residual improvement ratio relative to CSL for
    the restriction/prolongation diagnostics. Source plot:
    \texttt{restriction\_tup\_failure\_diagnostics/120\_residual\_improvement\_ratio\_vs\_csl.png}.}
    \label{fig:results-1d-restriction-improvement}
\end{figure}

\begin{figure}[htbp]
    \centering
    \includegraphics[width=\linewidth]{figures/ch6/121_gate_kept_frequency_exact_vs_learned.png}
    \caption{Modal gate acceptance for exact-low transfer
    versus learned restriction. Source plot:
    \texttt{restriction\_tup\_failure\_diagnostics/121\_gate\_kept\_frequency\_exact\_vs\_learned.png}.
    The exact-low plus $T_\uparrow$ path is accepted in only a small fraction
    of modes, but those modes carry enough solver value to reduce FGMRES
    iterations.}
    \label{fig:results-1d-gate-kept}
\end{figure}

The one-dimensional conclusion is therefore precise. A solution-transfer U-Net
does learn meaningful information about the high-frequency solution. Used
raw, that information can produce a worse residual than the zero start. Used
selectively, especially inside a preconditioned residual-correction pathway,
it can reduce FGMRES iterations. The scientific lesson is not that a neural
network replaces GMRES. It is that field prediction must be translated into
residual-compatible corrections before it becomes a solver tool.

% ============================================================
\section{Two-dimensional FD/PML results}
\label{sec:results-2d-fdpml}
% ============================================================

The two-dimensional experiments test the same principle in the realistic
finite-difference/PML setting. The grid is $512\times512$ with a PML of width
112 on each side, leaving a $288\times288$ physical interior. The frequency
pairs are
\begin{equation}
    16\to32,
    \qquad
    32\to64,
    \qquad
    64\to128.
\end{equation}
For each pair, the dataset contains $9600$ full-grid FD/PML samples with
complex Gaussian sources. The imaginary source channel is stored, and the
dataset audits passed. This matters because the two-dimensional question is
not merely whether the network predicts the physical interior. It is whether
the full FD/PML output is compatible with the same discrete operator used by
the solver.

The main model is a full-grid flux-trained \texttt{TransferUNet} with
\texttt{base\_ch=32} and five levels. The input is the real and imaginary
parts of the low-frequency field; the target is the real and imaginary parts
of the high-frequency field. The loss is the full-grid complex relative
$\ell^2$ loss. Unlike the older interior-focused models, the PML region is
included in training.

% ------------------------------------------------------------
\subsection{Full-grid field prediction}
\label{subsec:results-2d-field}
% ------------------------------------------------------------

Table~\ref{tab:results-2d-validation} reports the best validation errors. The
field errors are larger than in the one-dimensional Dirichlet study, as
expected: the two-dimensional problem includes a full FD/PML operator, a
larger grid, and more complex source fields. Nevertheless, the models learn
nontrivial transfer maps for all three pairs. The interior errors are
consistently lower than the full-grid errors, indicating that the PML remains
the harder region to model.

\begin{table}[htbp]
\centering
\caption{Best validation errors for the full-grid FD/PML 2D
\texttt{TransferUNet} models trained with complex relative $\ell^2$ loss.}
\label{tab:results-2d-validation}
\begin{tabular}{lcc}
\toprule
Pair & Full-grid validation RelL2 & Interior validation RelL2 \\
\midrule
$16\to32$ & $0.240$ & $0.190$ \\
$32\to64$ & $0.321$ & $0.252$ \\
$64\to128$ & $0.339$ & $0.276$ \\
\bottomrule
\end{tabular}
\end{table}

\begin{figure}[htbp]
    \centering
    \includegraphics[width=\linewidth]{figures/ch6/training_64_128.png}
    \caption{Training and validation curves for the full-grid FD/PML
    \texttt{TransferUNet} on the $64\to128$ pair. Source plot:
    \texttt{experiments/2d/presentation\_plots/flux\_full\_N9600/training\_64\_128.png}.
    The curve shows the hardest tested pair under the full-grid complex
    relative $\ell^2$ objective, with both full-grid and interior validation
    errors monitored during training.}
    \label{fig:results-2d-training-64-128}
\end{figure}

The result is a useful but incomplete success. Field prediction alone is not
the decisive metric, but it establishes that the full-grid models have learned
substantial information about $u_H$ from $u_L$. The next question is whether
that information survives application of the high-frequency FD/PML operator.

% ------------------------------------------------------------
\subsection{Warm-start FGMRES with true residual evaluation}
\label{subsec:results-2d-solver}
% ------------------------------------------------------------

The solver evaluation uses true-residual FGMRES with a CSL preconditioner at
$\beta=0.1$. For each frequency pair, $10$ test samples are evaluated with a
maximum budget of $K=30$ iterations. The CSL preconditioner is applied through
an exact sparse LU solve of the shifted operator. The reported initial
residual is
\begin{equation}
    r_0^{\mathrm{rel}}
    =
    \frac{\|b-A_Hx_0\|_2}{\|b\|_2}.
\end{equation}
This is the correct residual-facing metric for judging whether a warm start is
better than the cold initial guess.

\begin{table}[htbp]
\centering
\caption{Two-dimensional FD/PML warm-start FGMRES results with true-residual
evaluation, CSL $\beta=0.1$, $10$ test samples, and budget $K=30$. The
\texttt{flux\_full\_raw} model improves field error and iteration counts
relative to the cold start, but its initial residual remains above the
cold-start value.}
\label{tab:results-2d-solver}
\begin{tabular}{llccc}
\toprule
Pair & Method & Full error & $\|r_0\|_2/\|b\|_2$ & Mean iterations \\
\midrule
$16\to32$ & Cold & $1.000$ & $1.00$ & $10.0$ \\
          & Depth5 zero & $0.570$ & $15.76$ & $9.1$ \\
          & Flux-full raw & $0.238$ & $1.49$ & $9.0$ \\
\midrule
$32\to64$ & Cold & $1.000$ & $1.00$ & $14.0$ \\
          & Depth5 zero & $0.590$ & $11.25$ & $13.5$ \\
          & Flux-full raw & $0.352$ & $1.93$ & $13.0$ \\
\midrule
$64\to128$ & Cold & $1.000$ & $1.00$ & $28.9$ \\
           & Depth5 zero & $0.616$ & $6.40$ & $28.0$ \\
           & Flux-full raw & $0.276$ & $1.86$ & $27.0$ \\
\bottomrule
\end{tabular}
\end{table}

The main result is that \texttt{flux\_full\_raw} is the best current 2D learned
warm start across all three tested pairs. It gives the lowest full-grid field
error among the listed methods and reduces the mean FGMRES iteration count
relative to the cold start:
\[
    10.0 \to 9.0,
    \qquad
    14.0 \to 13.0,
    \qquad
    28.9 \to 27.0.
\]
These are modest but consistent improvements.

The caution is equally important. The initial residual remains worse than the
cold start for every pair: $1.49$, $1.93$, and $1.86$, compared with $1.00$
for $x_0=0$. Therefore the correct claim is not that the neural warm start
already solves the residual problem. The correct claim is that full-grid
FD/PML training moves the neural warm start much closer to the cold-residual
regime while improving FGMRES iteration counts across the tested pairs.

\begin{figure}[htbp]
    \centering
    \includegraphics[width=\linewidth]{figures/ch6/gmres_clean_64_128.png}
    \caption{Clean true-residual FGMRES convergence curves for the
    $64\to128$ FD/PML pair. Source plot:
    \texttt{experiments/2d/presentation\_plots/flux\_full\_solver\_eval\_beta0p3\_clean/gmres\_clean\_64\_128.png}.
    This is the hardest tested transfer pair, so the plot gives the clearest
    solver-facing view of the cold start and neural warm-start behaviour.}
    \label{fig:results-2d-gmres-clean-64-128}
\end{figure}

\begin{figure}[htbp]
    \centering
    \includegraphics[width=\linewidth]{figures/ch6/2d_initial_residual_bars.png}
    \caption{Initial true residual
    $\|b-A_Hx_0\|_2/\|b\|_2$ for cold, old neural variants, and
    \texttt{flux\_full\_raw}. This figure should emphasise the central
    caveat: full-grid FD/PML training greatly reduces the residual pathology
    of older models, but the raw neural warm start still begins above the
    cold-start residual.}
    \label{fig:results-2d-r0-bars}
\end{figure}

The depth5-zero comparison illustrates why PML behaviour is solver-facing, not
cosmetic. Zeroing the PML can reduce full-grid field error relative to a bad
raw PML output, but it may create a very large initial residual. In
Table~\ref{tab:results-2d-solver}, the depth5-zero residuals are
$15.76$, $11.25$, and $6.40$. This reproduces the one-dimensional lesson in a
different form: a field-level operation that looks helpful can be harmful
after applying the operator.

% ------------------------------------------------------------
\subsection{PML compatibility}
\label{subsec:results-2d-pml}
% ------------------------------------------------------------

The older 2D models had good interior field behaviour but unreliable PML
outputs. The full-grid FD/PML training changes this. The raw PML output of
\texttt{flux\_full\_raw} is now useful enough that keeping it is better than
zeroing it. This is the main practical improvement of the new 2D pipeline.

\begin{figure}[htbp]
    \centering
    \includegraphics[width=\linewidth]{figures/ch6/2d_pml_energy_all_pairs.png}
    \caption{PML energy ratios for old depth5 variants,
    PML-zeroed variants, and \texttt{flux\_full\_raw}. The intended message is
    that full-grid FD/PML training makes the raw PML output solver-compatible:
    the PML is no longer a region to remove by default.}
    \label{fig:results-2d-pml-energy}
\end{figure}

This result should be interpreted carefully. The PML-compatible model still
does not beat the cold start in initial true residual. It does, however, avoid
the extreme residual blow-up caused by earlier PML handling and gives
consistent iteration-count reductions. In the language of the one-dimensional
study, full-grid FD/PML training acts like a crude compatibility constraint:
it does not provide exact spectral gating, but it prevents the learned field
from violating the discrete operator as severely in the absorbing layer.

\begin{table}[htbp]
\centering
\caption{Recommended compact table for the final thesis layout. It combines
field accuracy, PML behaviour, residual compatibility, and iteration count.
Only the columns needed for solver interpretation should be retained in the
final version.}
\label{tab:results-2d-compact}
\begin{tabular}{llcccc}
\toprule
Pair & Method & Full error & PML ratio & $\|r_0\|_2/\|b\|_2$ & Mean iterations \\
\midrule
$16\to32$ & Cold & $1.000$ & -- & $1.00$ & $10.0$ \\
          & Depth5 zero & $0.570$ & $0.000$ & $15.76$ & $9.1$ \\
          & Flux-full raw & $0.238$ & -- & $1.49$ & $9.0$ \\
\midrule
$32\to64$ & Cold & $1.000$ & -- & $1.00$ & $14.0$ \\
          & Depth5 zero & $0.590$ & $0.000$ & $11.25$ & $13.5$ \\
          & Flux-full raw & $0.352$ & -- & $1.93$ & $13.0$ \\
\midrule
$64\to128$ & Cold & $1.000$ & -- & $1.00$ & $28.9$ \\
           & Depth5 zero & $0.616$ & $0.000$ & $6.40$ & $28.0$ \\
           & Flux-full raw & $0.276$ & -- & $1.86$ & $27.0$ \\
\bottomrule
\end{tabular}
\end{table}

The table includes a PML-ratio column because the final plotted diagnostics
should report it, but the present draft does not insert unlisted numerical
PML-ratio values. This avoids giving false precision before the final figure
and table pass.

% ------------------------------------------------------------
\subsection{Interpretation of the 2D results}
\label{subsec:results-2d-interpretation}
% ------------------------------------------------------------

The two-dimensional results transfer the one-dimensional lesson, but only in a
weaker and more realistic form. In 1D, analytical eigenpairs allow exact modal
gating. In 2D FD/PML, the operator is large, complex-valued, and non-normal.
There is no equally simple modal decomposition in the experiments reported
here. Therefore the residual and FGMRES curves become the decisive diagnostics.

The result is encouraging because the full-grid FD/PML model improves
iterations for all tested pairs and avoids the severe PML-induced residual
pathology of older models. It is also incomplete because the raw neural warm
start still begins with a true residual above the cold-start residual. The
next solver-facing step is therefore not simply to train a lower field-loss
model. The next step is to incorporate residual-compatible selection, loss
terms, or gating into the 2D method, following the one-dimensional evidence
that only selected learned components improve the preconditioned solve.

% ============================================================
\section{Chapter summary}
\label{sec:results-summary}
% ============================================================

The results support three conclusions.

First, field accuracy is not solver accuracy. In the 1D Dirichlet problem, a
U-Net with field error near $0.06$ can produce an initial residual roughly
$11.6$ times larger than $\|b\|_2$. The analytical eigenbasis shows exactly
why: $c_k(r_0)=\lambda_k c_k(e_0)$.

Second, learned transfer is still useful when made residual-aware. Spectral
filtering and residual gating reduce the harmful high-$|\lambda|$ components,
and an exact-low plus gated-$T_\uparrow$ preconditioner reduces the 1D FGMRES
iteration count from about $16.0$ to about $12.65$. The successful learned
component is selective, not wholesale.

Third, the 2D FD/PML experiments show cautious solver-facing progress. Full-grid
FD/PML training makes the raw PML output useful, brings the learned warm start
much closer to the cold residual regime, and improves FGMRES iteration counts
on $16\to32$, $32\to64$, and $64\to128$. However, the initial true residual
remains above the cold-start value, so the method should not be presented as a
finished neural accelerator. It is best understood as evidence that learned
frequency transfer can provide solver-relevant information, provided that the
field prediction is constrained or gated by residual metrics.
